\section*{Resumen}
\addcontentsline{toc}{section}{Resumen}

Este informe presenta el primer entregable del Trabajo Parcial de Programación Concurrente y
Distribuida. El grupo propone aplicar K-means a los registros de viajes del taxi amarillo de Nueva
York \parencite{nyctlc} para caracterizar patrones de movilidad urbana, un insumo para la
planificación del transporte que persigue el Objetivo de Desarrollo Sostenible~11. El algoritmo se
implementará en PC2 en Go, sin librerías de terceros, con goroutines, \texttt{sync.WaitGroup} y
\texttt{sync.Mutex}.

La revisión bibliográfica reúne nueve artículos publicados entre 2021 y 2026, tres por integrante.
Se organizan en tres temas: paralelismo en memoria compartida, sincronización en entornos
distribuidos y reducción del trabajo mediante poda o aproximación. Cada artículo tiene una ficha
con los campos del enunciado, y cada tema cierra con una comparación que orienta el diseño del
algoritmo concurrente.

El dataset corresponde a enero de 2024. Un pipeline reproducible aplica siete reglas de limpieza,
las contrasta con una auditoría independiente en SQL y conserva el 95,5\,\% de los viajes, muy por
encima del millón de registros que exige el enunciado.

\textit{Palabras clave:} K-means, programación concurrente, Go, NYC TLC, movilidad urbana, ODS 11.
